\documentclass[11pt,a4paper]{article}
\usepackage[utf8]{inputenc}
\usepackage[vietnamese]{babel}
\usepackage{amsmath,amssymb}
\usepackage{geometry}
\geometry{a4paper, total={180mm,260mm}, left=15mm, right=15mm, top=18mm, bottom=18mm}

\usepackage{tikz}
\usetikzlibrary{shapes.geometric, arrows.meta, positioning, calc}
\tikzset{arr/.style={-{Stealth[scale=0.8]}}}
\usepackage{tcolorbox}
\tcbuselibrary{skins,breakable}
\usepackage{fontawesome}
\usepackage{fancyhdr}
\usepackage{enumitem}

\pagestyle{fancy}
\fancyhf{}
\renewcommand{\headrulewidth}{0.5pt}
\renewcommand{\footrulewidth}{0.5pt}
\fancyhead[L]{\small\bfseries\color{violet!85!black}\faBook\ SỔ TAY KIẾN THỨC TOÁN 2}
\fancyhead[R]{\small\bfseries\color{violet!85!black}Lớp Toán Thầy Lĩnh -- loptoanthaylinh.vn}
\fancyfoot[C]{\small\color{gray}-- Sổ tay kiến thức: Bài 07. Phép cộng có tổng là số tròn chục --}

\setlength{\parindent}{0pt}
\setlength{\parskip}{0pt}

\begin{document}

% --- BANNER BÀI HỌC ---
\begin{tcolorbox}[
  enhanced,
  colback=violet!8!white,
  colframe=violet!70!black,
  arc=3mm, boxrule=1.2pt,
  left=3mm, right=3mm, top=2mm, bottom=2mm,
  halign=center
]
  {\LARGE\bfseries\color{violet!90!black}\faBookmark\ BÀI 07: PHÉP CỘNG CÓ TỔNG LÀ SỐ TRÒN CHỤC}\\[1mm]
  {\itshape\footnotesize\color{black!75} Kĩ thuật cộng qua 10 làm tròn chục dạng 26 + 4 và 36 + 24 trên 1 trang A4}
\end{tcolorbox}

\vspace{-2mm}

\noindent
\begin{minipage}[t]{0.49\textwidth}\vspace{0pt}
\begin{tcolorbox}[
  enhanced,
  colback=violet!4!white,
  colframe=violet!75!black,
  arc=2mm, boxrule=0.8pt,
  left=2mm, right=2mm, top=1.8mm, bottom=1.8mm,
  title={\footnotesize\bfseries\color{white}\faKey\ 1. Dạng 1: Số có 2 chữ số cộng số có 1 chữ số}
]
{\footnotesize
\begin{itemize}\setlength{\itemsep}{1pt}\setlength{\leftskip}{-2.5mm}
  \item \textbf{Ví dụ:} $26 + 4 = 30$.
  \item \textbf{Đặt tính thẳng cột:}\\
  Viết $4$ thẳng cột với $6$ (hàng đơn vị).
  \item \textbf{Quy trình tính:}\\
  - Hàng đơn vị: $6 + 4 = 10$, \textbf{viết 0, nhớ 1}.\\
  - Hàng chục: $2$ \textbf{thêm 1 bằng 3}, viết $3$.\\
  $\implies 26 + 4 = \mathbf{30}$ (số tròn chục).
  \item \textbf{Các cặp cộng ra 10:}\\
  $9+1=10$; $8+2=10$; $7+3=10$; $6+4=10$; $5+5=10$.
\end{itemize}
}
\end{tcolorbox}
\end{minipage}
\hfill
\begin{minipage}[t]{0.49\textwidth}\vspace{0pt}
\begin{tcolorbox}[
  enhanced,
  colback=teal!4!white,
  colframe=teal!75!black,
  arc=2mm, boxrule=0.8pt,
  left=2mm, right=2mm, top=1.8mm, bottom=1.8mm,
  title={\footnotesize\bfseries\color{white}\faSliders\ 2. Dạng 2: Số có 2 chữ số cộng số có 2 chữ số}
]
{\footnotesize
\begin{itemize}\setlength{\itemsep}{1pt}\setlength{\leftskip}{-2.5mm}
  \item \textbf{Ví dụ:} $36 + 24 = 60$.
  \item \textbf{Đặt tính thẳng cột:}\\
  Đơn vị thẳng đơn vị, chục thẳng chục.
  \item \textbf{Quy trình tính:}\\
  - Hàng đơn vị: $6 + 4 = 10$, \textbf{viết 0, nhớ 1}.\\
  - Hàng chục: $3 + 2 = 5$, \textbf{thêm 1 bằng 6}, viết $6$.\\
  $\implies 36 + 24 = \mathbf{60}$.
  \item \textbf{Kĩ năng nhớ:}\\
  Luôn nhớ $1$ vào tổng hàng chục sau khi cộng.
\end{itemize}
}
\end{tcolorbox}
\end{minipage}

\vspace{-1.5mm}
\begin{tcolorbox}[
  enhanced,
  colback=orange!4!white,
  colframe=orange!80!black,
  arc=2mm, boxrule=0.8pt,
  left=2mm, right=2mm, top=1.5mm, bottom=1.5mm,
  title={\footnotesize\bfseries\color{white}\faSitemap\ 3. Sơ đồ tư duy cơ chế nhớ 1 sang hàng chục}
]
\centering
\begin{tikzpicture}[x=1cm, y=1cm]
  \node[rectangle, rounded corners=2pt, draw=violet!80!black, fill=violet!10, text width=3.8cm, align=center, font=\scriptsize\bfseries] at (0,0) {
    HÀNG ĐƠN VỊ\\$6 + 4 = 10$\\Viết $\mathbf{0}$ ở hàng đơn vị
  };
  \draw[->, very thick, red!85!black] (2,0) -- node[above, font=\tiny\bfseries\color{red!90!black}] {Nhớ 1 chục} (3.5,0);
  \node[rectangle, rounded corners=2pt, draw=teal!80!black, fill=teal!10, text width=4.5cm, align=center, font=\scriptsize\bfseries] at (5.8,0) {
    HÀNG CHỤC\\$3 + 2 = 5$\\Thêm $1$ nhớ: $5 + 1 = \mathbf{6}$
  };
\end{tikzpicture}
\end{tcolorbox}

\vspace{-1.5mm}
\noindent
\begin{minipage}[t]{0.49\textwidth}\vspace{0pt}
\begin{tcolorbox}[
  enhanced,
  colback=green!4!white,
  colframe=green!60!black,
  arc=2mm, boxrule=0.8pt,
  left=2mm, right=2mm, top=1.8mm, bottom=1.8mm,
  title={\footnotesize\bfseries\color{white}\faCheckCircle\ 4. Ví dụ thực hành chuẩn mực}
]
{\footnotesize
\textbf{1. Đặt tính rồi tính:}\\
\begin{tabular}{r@{\,}lr@{\,}l}
  $4$ & $7$ & $5$ & $3$ \\
$+$ & & $+$ & \\
    & $3$ & $2$ & $7$ \\ \hline
  $\mathbf{5}$ & $\mathbf{0}$ & $\mathbf{8}$ & $\mathbf{0}$
\end{tabular}
\quad
\begin{tabular}{l}
  $\bullet$ $7+3=10$, viết $0$ nhớ $1$.\\
  $\bullet$ $4$ thêm $1$ là $5$, viết $5$.\\
  $\bullet$ $3+7=10$, viết $0$ nhớ $1$.\\
  $\bullet$ $5+2=7$, thêm $1$ là $8$.
\end{tabular}
\textbf{2. Tính nhẩm:}\\
$75 + 5 = 80$; $48 + 12 = 60$.
}
\end{tcolorbox}
\end{minipage}
\hfill
\begin{minipage}[t]{0.49\textwidth}\vspace{0pt}
\begin{tcolorbox}[
  enhanced,
  colback=red!4!white,
  colframe=red!75!black,
  arc=2mm, boxrule=0.8pt,
  left=2mm, right=2mm, top=1.8mm, bottom=1.8mm,
  title={\footnotesize\bfseries\color{white}\faExclamationTriangle\ 5. Mẹo tránh bẫy trắc nghiệm}
]
{\footnotesize
\begin{itemize}\setlength{\itemsep}{1pt}\setlength{\leftskip}{-2.5mm}
  \item \textbf{Bẫy quên cộng 1 nhớ sang hàng chục:}\\
  Nhiều bạn tính $26 + 4 = 20$ (vì $6+4=10$ viết 0 rồi hạ 2 xuống!). Đây là lỗi sai rất nặng. Phải \textbf{thêm 1} vào hàng chục: $2 + 1 = 3 \implies 30$.
  \item \textbf{Bẫy đặt tính:}\\
  Số có 1 chữ số phải viết thẳng cột hàng đơn vị (bên phải).
  \item Kết quả luôn tận cùng bằng chữ số $0$.
\end{itemize}
}
\end{tcolorbox}
\end{minipage}


\end{document}
